% Options for packages loaded elsewhere
\PassOptionsToPackage{unicode}{hyperref}
\PassOptionsToPackage{hyphens}{url}
\documentclass[
  ignorenonframetext,
]{beamer}
\newif\ifbibliography
\usepackage{pgfpages}
\setbeamertemplate{caption}[numbered]
\setbeamertemplate{caption label separator}{: }
\setbeamercolor{caption name}{fg=normal text.fg}
\beamertemplatenavigationsymbolsempty
% remove section numbering
\setbeamertemplate{part page}{
  \centering
  \begin{beamercolorbox}[sep=16pt,center]{part title}
    \usebeamerfont{part title}\insertpart\par
  \end{beamercolorbox}
}
\setbeamertemplate{section page}{
  \centering
  \begin{beamercolorbox}[sep=12pt,center]{section title}
    \usebeamerfont{section title}\insertsection\par
  \end{beamercolorbox}
}
\setbeamertemplate{subsection page}{
  \centering
  \begin{beamercolorbox}[sep=8pt,center]{subsection title}
    \usebeamerfont{subsection title}\insertsubsection\par
  \end{beamercolorbox}
}
% Prevent slide breaks in the middle of a paragraph
\widowpenalties 1 10000
\raggedbottom
\AtBeginPart{
  \frame{\partpage}
}
\AtBeginSection{
  \ifbibliography
  \else
    \frame{\sectionpage}
  \fi
}
\AtBeginSubsection{
  \frame{\subsectionpage}
}
\usepackage{iftex}
\ifPDFTeX
  \usepackage[T1]{fontenc}
  \usepackage[utf8]{inputenc}
  \usepackage{textcomp} % provide euro and other symbols
\else % if luatex or xetex
  \usepackage{unicode-math} % this also loads fontspec
  \defaultfontfeatures{Scale=MatchLowercase}
  \defaultfontfeatures[\rmfamily]{Ligatures=TeX,Scale=1}
\fi
\usepackage{lmodern}
\ifPDFTeX\else
  % xetex/luatex font selection
\fi
% Use upquote if available, for straight quotes in verbatim environments
\IfFileExists{upquote.sty}{\usepackage{upquote}}{}
\IfFileExists{microtype.sty}{% use microtype if available
  \usepackage[]{microtype}
  \UseMicrotypeSet[protrusion]{basicmath} % disable protrusion for tt fonts
}{}
\makeatletter
\@ifundefined{KOMAClassName}{% if non-KOMA class
  \IfFileExists{parskip.sty}{%
    \usepackage{parskip}
  }{% else
    \setlength{\parindent}{0pt}
    \setlength{\parskip}{6pt plus 2pt minus 1pt}}
}{% if KOMA class
  \KOMAoptions{parskip=half}}
\makeatother
\setlength{\emergencystretch}{3em} % prevent overfull lines
\providecommand{\tightlist}{%
  \setlength{\itemsep}{0pt}\setlength{\parskip}{0pt}}
\usepackage{bookmark}
\IfFileExists{xurl.sty}{\usepackage{xurl}}{} % add URL line breaks if available
\urlstyle{same}
\hypersetup{
  hidelinks,
  pdfcreator={LaTeX via pandoc}}

\author{\texorpdfstring{}{}}
\date{}

\begin{document}

\section{World models, active inference, boundaries, and forward
compatibility}\label{world-models-active-inference-boundaries-and-forward-compatibility}

\begin{frame}{World models from code}
\protect\phantomsection\label{world-models-from-code}
The central theoretical claim of COGANT is that source code implicitly
defines a generative model of system behaviour, a claim that is
structurally analogous to the \emph{world-model} line of
reinforcement-learning work exemplified by Dreamer V3
{[}@hafner2023dreamerv3{]}: an encoder maps observations to a latent
state, latent dynamics predict forward, and a decoder maps back. The
Dreamer architecture learns the world model from observation
trajectories. COGANT extracts it symbolically from the program graph;
the comparison is productive precisely because it clarifies that COGANT
produces an explicit, interpretable state-space and transition structure
where Dreamer produces an opaque learned latent. Both are generative
models that a downstream Active Inference agent {[}@friston2010free;
@parr2022active; @dacosta2020active; @heins2022pymdp{]} can treat as the
\(p(o, s)\) component of a variational free-energy functional.
\end{frame}

\begin{frame}[fragile]{Active inference and program behavior}
\protect\phantomsection\label{active-inference-and-program-behavior}
The state-space IR in COGANT's pipeline (states, actions, transitions,
observations) shares structural parallels with \textbf{active inference}
formulations {[}@friston2010free; @parr2022active{]}, where an agent
maintains beliefs about hidden states and selects actions to minimize
prediction error. The discrete-state synthesis presented in
{[}@dacosta2020active{]} is the closest formal target of COGANT's
compilation: variables, actions, observation modalities, and transition
structures in the Generalized Notation Notation bundle map directly onto
the tuples required by a discrete-state active inference agent, and the
step-by-step construction protocol of {[}@smith2022stepbystep{]} can be
followed literally against those bundles. PyMDP {[}@heins2022pymdp{]}
provides a reference Python runtime that executes exactly this form of
agent, making it a natural downstream consumer of COGANT exports. In the
program analysis context, the ``agent'' is the analysis pipeline itself:
it observes code artifacts, maintains beliefs about program behavior
(the state-space model), and refines those beliefs as new evidence
(dynamic traces, coverage data) arrives.

This connection is analogical: the \texttt{ConfidenceModel} in
\texttt{../cogant/py/cogant/translate/confidence.py} aggregates evidence
and penalties in a way that suggests belief revision, but it is not a
Bayesian posterior. Future work could formalize a tighter link by
casting rule application as variational inference, where a fixpoint
would represent an approximate posterior over program semantics.
\end{frame}

\begin{frame}[fragile]{Boundaries}
\protect\phantomsection\label{boundaries}
COGANT does not subsume formal verification, interactive theorem
proving, or full interprocedural pointer analysis unless implemented as
explicit future stages.
\href{../cogant/docs/reference/implementation_status.md}{\texttt{../cogant/docs/reference/implementation\_status.md}}
marks Rust acceleration and additional parsers as staged; the manuscript
should be read together with that table for up-to-date scope.
\end{frame}

\begin{frame}[fragile]{Forward compatibility}
\protect\phantomsection\label{forward-compatibility}
Promoting COGANT into
\href{../../../projects/}{\texttt{../../../projects/}} integrates
manuscript PDF rendering with the template's validation gates.
Cross-references in this folder use paths \textbf{relative to these
Markdown files} (for example
\href{../cogant/docs/}{\texttt{../cogant/docs/}}) so links stay stable
when the tree moves.
\end{frame}

\end{document}
